\section{Conclusão}

Após a aplicação dos modelos para tentar prever o escore de cada aluno, ficou evidenciado que a rede neural seria a melhor escolha para este conjunto de dados, com erros baixos em relação ao real valor dos escores, o estudo evidencia que é viável tentar prever a qualidade no desempenho escolar.


O bom resultado do ajuste serve como munição para encorajar trabalhos futuros que busquem ter acesso a dados reais relacionados a desempenho estudantil. Assim buscando entender se o comportamento de dados reais é similar ao simulado.

O trabalho consegue atingir seu objetivo de estudar o conjunto de dados e aplicar os modelos para entender qual seria a melhor opção para prever o escore dos alunos, apesar de se tratarem de dados fictícios, o resultado é muito positivo e chama atenção para opções futuras implementadas na sociedade.
